\hmsubsection{Communication Between Systems}{FS}{}\label{sec:comm}
The Raspberry Pi 5 and the OpenRB-150 communicate over USB serial. The
link runs at 115\,200~bps, 8N1. On the Pi side, it appears as the
\texttt{/dev/ttyACM0} device. On the OpenRB-150 side, it uses the
on-board USB serial port.
USB was selected instead of a hardware UART on the Pi for two reasons.
First, the USB cable also powers the OpenRB-150 logic circuit, reducing
the required wiring. Second, the OpenRB-150 has a front-facing USB-C
connector, which is easier to access than the GPIO pins on the Pi. The
12 V motor power is delivered separately to the OpenRB-150's motor power
input, as described in Section \ref{sec:hardware-integration}.
\subsubsection*{Message format}
Each message consists of one line of ASCII terminated by a newline. The
host sends commands, and the firmware sends status replies.
\subsubsection*{Host side}
The host-side implementation is contained in the \texttt{SerialBridge}
class. There is one method per command. Each method writes the command
line and then reads the reply.
\texttt{MOVE} blocks until \texttt{READY} or a timeout. \texttt{LIMIT}
and \texttt{SCAN} return on the first \texttt{ACK}. The default cycle
timeout is 10 s, which is about twice the longest motion time we
recorded in testing. If the timeout fires before \texttt{READY},
\texttt{SerialBridge} raises a typed exception. The supervisor handles
the exception by aborting the cycle and sending the arm back to the
scan pose.
\subsubsection*{Firmware side}
The firmware reads bytes from USB serial in the main loop and parses
each newline-terminated line into a command and zero or more integer
arguments. \texttt{MOVE} commands are checked against the joint limits
in Section \ref{sec:joint-limits}. If the check passes, the firmware
writes all five goal positions in one \texttt{syncWrite()} call. The
firmware then polls the present positions of all five motors and
sends \texttt{BUSY} while any motor is still moving. Once every motor
is within tolerance of its goal, the firmware sends \texttt{READY}.
\subsubsection*{Why text instead of CRC framing}
Text framing was selected instead of a binary frame with CRC for three
reasons.
First, the protocol can be tested from a serial terminal. During
bring-up, commands were entered manually to verify firmware behaviour. A
binary protocol would have required a separate Python tool from the
start of development.
Second, the cable between the Pi and the OpenRB-150 is shorter than
one metre, and the laboratory environment is electrically quiet. The
available development observations did not indicate transmission errors
in this bench setup. A CRC check was therefore not expected to provide
measurable benefit in the current configuration.
Third, the firmware is limited in scope. A binary parser would have
added complexity without meaningful space savings.
Sortify \cite{sortify} used a binary frame with CRC-16 because they
intended to place the arm on a moving vehicle. They also reported no CRC
errors in their stress test on a bench setup comparable to this one.
That result is consistent with the observations in this project and
supports the current choice of text framing. If the cell is later placed
on Fenris, where motor noise and cable routing may change, the wire
protocol can be switched to binary without changing anything above the
\texttt{SerialBridge} layer.
